% sup4_02_zero_b: Supplement part 4, zero-field kernel bounds at every point, II: the free massive propagator, lattice
% convolution, and the remainder G - C^xi (the gain of the resolvent comparison). Input through INPUTS_part4.tex.

\subsection{Zero-field kernel bounds at every point of a box, II: the free propagator and the remainder $\breve G-C^\xi$}\label{sup4:ssec:zeroB}

\begin{lemma}[free massive propagator smeared over a block]\label{sup4:lem:C1}
Let $\xi=L^{-J}$, $J\ge0$, $L$ odd, $\mu^2\in[\frac12,2]$, and $C_\mu\defeq(-\Delta^\xi+\mu^2)^{-1}$ on $\xi\Z^d$. Put $\delta_C\defeq(88d)^{-1/2}$. Then for every unit block $B(y)$, every $x$, $\tilde x$ with $|x-\tilde x|\le1$ and every direction,
\[
|(C_\mu\mathbf 1_{B(y)})(x)|,\ |D(C_\mu\mathbf 1_{B(y)})(x)|\le c_Ce^{-\delta_C|x-y|},\qquad |x-\tilde x|^{-\alpha}|D(C_\mu\mathbf 1_{B(y)})(\tilde x)-D(C_\mu\mathbf 1_{B(y)})(x)|\le c_C(\alpha)e^{-\delta_C\min(|x-y|,|\tilde x-y|)},
\]
with $c_C$, $c_C(\alpha)$ depending only on $(d,\alpha)$.
\end{lemma}

\begin{proof}
Let $N\defeq1/\xi=L^J$ (odd) and $B(y)=y+([0,1)\cap\xi\Z)^d$. With $\hat f(p)\defeq\sum_x\xi^df(x)e^{-ip\cdot x}$ on $B_\xi\defeq[-\pi/\xi,\pi/\xi]^d$, the geometric sum $\sum_{n=0}^{N-1}\xi e^{-it\xi n}=e^{-i(1-\xi)t/2}g(t)$, $g(t)\defeq\xi\sin(t/2)/\sin(\xi t/2)$, gives $\hat{\mathbf 1}_{B(y)}(p)=e^{-ip\cdot y}u(p)$ with $u(p)\defeq e^{-i(1-\xi)(p_1+\dots+p_d)/2}\prod_\nu g(p_\nu)$. Hence, for $M\in\{1,\hat\partial_\nu,\hat\partial_\nu(e^{ip\cdot v'}-1)/|v'|^\alpha\}$ ($\hat\partial_\nu(p)=(e^{i\xi p_\nu}-1)/\xi$, $v'\defeq\tilde x-x\in\xi\Z^d$, $0<|v'|\le1$), the three quantities of the statement are
\[
(2\pi)^{-d}\int_{B_\xi}e^{ip\cdot(x-y)}\,\frac{M(p)\,u(p)}{\Delta^\xi(p)+\mu^2}\,dp,\qquad \Delta^\xi(p)=\sum_\nu\frac4{\xi^2}\sin^2\frac{\xi p_\nu}2.
\]
\emph{Step 1 (periodicity).} Every factor is $(2\pi/\xi)$-periodic in each $p_\nu$: $x-y$, $v'\in\xi\Z^d$; $\Delta^\xi$ and $\hat\partial$ are trigonometric in $\xi p$; and $u$ is periodic because $N$ is odd: $g(t+2\pi/\xi)=\xi(-1)^N\sin(t/2)/(-\sin(\xi t/2))=g(t)$, and the phase gains $e^{-i\pi(N-1)}=1$. The integrand is entire in each $p_\nu$ ($g$ is entire).

\emph{Step 2 (the denominator on the strip).} Let $p=q+i\sigma$, $q\in B_\xi$, $|\sigma|_\infty\le\delta_C\le\frac14$. From $\sin^2(a+ib)-\sin^2a=\sinh^2b\,(\sin^2a-\cos^2a)+\frac i2\sin2a\,\sinh2b$ with $a=\xi q_\nu/2$, $b=\xi\sigma_\nu/2$, and $\sinh x\le1.13x$, $\sinh2x\le3.1x$ for $0\le x\le\frac12$: $|\Delta^\xi_\nu(p)-\Delta^\xi_\nu(q)|\le1.28\sigma_\nu^2+3.1|\sigma_\nu|\Delta^\xi_\nu(q)^{1/2}$. Summing over $\nu$ and using $\sum_\nu\Delta_\nu^{1/2}\le\sqrt d\,\Delta^{1/2}$ and $3.1\delta\sqrt d\,\Delta^{1/2}\le\frac14\Delta+9.7d\delta^2$: $|\Delta^\xi(p)-\Delta^\xi(q)|\le\frac14\Delta^\xi(q)+11d\delta_C^2\le\frac14\Delta^\xi(q)+\frac18\le\frac14(\Delta^\xi(q)+\mu^2)$. Hence $|\Delta^\xi(p)+\mu^2|\ge\frac34(\Delta^\xi(q)+\mu^2)\ge c\,(1+|q|)^2$, since $\Delta^\xi(q)\ge(4/\pi^2)|q|^2$ on $B_\xi$.

\emph{Step 3 (the numerator).} For $z=t+i\sigma$, $|t|\le\pi/\xi$, $|\sigma|\le\frac14$: $|\sin(z/2)|^2=\sin^2(t/2)+\sinh^2(\sigma/2)$ and $|\sin(\xi z/2)|^2=\sin^2(\xi t/2)+\sinh^2(\xi\sigma/2)\ge(\xi t/\pi)^2+(\xi\sigma/2)^2$. For $|t|\le\pi$ this gives $|g(z)|^2\le\max\{(\pi/2)^2,1.03\}\le2.5$; for $|t|>\pi$, $|\sin(z/2)|\le\cosh\frac18\le1.01$ gives $|g(z)|\le3.2/|t|$. In both cases $|g(z)|\le10/(1+|t|)$. The phase of $u$ has modulus $\le e^{d\delta_C/2}\le1.2$, $|\hat\partial_\nu(p)|\le1.3|p_\nu|\le1.3(1+|q|)$, and $|e^{ip\cdot v'}-1|/|v'|^\alpha\le C\min\{1,|q||v'|\}^\alpha|v'|^{-\alpha}+C|v'|^{1-\alpha}\le C(1+|q|)^\alpha$ on the strip. So $|M(p)|\le C(1+|q|)^{1+\alpha}$.

\emph{Step 4 (integrability).} By Steps 2--3 the modulus of the integrand on the shifted contour is $\le e^{-\sigma\cdot(x-y)}\,C(d)\,(1+|q|)^{\alpha-1}\prod_\nu(1+|q_\nu|)^{-1}\le e^{-\sigma\cdot(x-y)}C(d)\prod_\nu(1+|q_\nu|)^{-1-(1-\alpha)/d}$, using $(1+|q|)^d\ge\prod_\nu(1+|q_\nu|)$; this is integrable over $\R^d$ uniformly in $\xi$.

\emph{Step 5 (the shift).} Shift each $p_\nu$ to $p_\nu+i\delta_C\operatorname{sgn}(x-y)_\nu$; the two vertical sides of each rectangle cancel by Step 1, and the shifted integral carries $|e^{ip\cdot(x-y)}|=e^{-\delta_C|x-y|_1}$. For the H\"older quotient, $|x-y|\ge\min(|x-y|,|\tilde x-y|)-1$ costs a factor $e^{\delta_C}$.
\end{proof}

\begin{lemma}[the free propagator $C^\xi\defeq(-\Delta^\xi+1)^{-1}$ on $\xi\Z^d$, $\xi=L^{-J}$]\label{sup4:lem:C}
For all $S$, $T$: $|\Delta_TD_SC^\xi(x_1,x_2)|\le C_C'(\xi+\rho)^{-p(S,T)}e^{-\delta'\rho}\prod_{i\in T}|x_i-\tilde x_i|^\alpha$, with $C_C'$, $\delta'$ depending only on $(d,L,\alpha)$.
\end{lemma}

\begin{proof}
Let $G\defeq\mathsf T_J^{-1}$ of \cref{sup4:cor:Z1} on $\xi\Z^d$ ($\eta\to\xi$, $k\to J$) with $m^2=1$; its block term is $a_JQ^*Q$ on unit blocks. Then $(C^\xi)^{-1}=G^{-1}-a_JQ^*Q$, so $C^\xi=G+a_JGQ^*QC^\xi$, and the kernel of the second term is $a_J\sum_y(G\mathbf 1_{B(y)})(x_1)(C^\xi\mathbf 1_{B(y)})(x_2)$, since $(QC^\xi f)(y)=\ip{C^\xi\mathbf 1_{B(y)}}f$. $G$ obeys \eqref{sup4:eq:Z1}. By (A5) for $G\mathbf 1_{B(y)}$ and \cref{sup4:lem:C1} for $C^\xi\mathbf 1_{B(y)}$, and the unit-lattice sum as in \cref{sup4:lem:Z}, Step 3, the second term obeys $|\Delta_TD_S\cdot|\le Ce^{-\delta''\rho}\prod|\cdot|^\alpha\le C'(\xi+\rho)^{-p}e^{-\delta''\rho/2}\prod|\cdot|^\alpha$.
\end{proof}

\begin{lemma}[lattice convolution, $d=3$]\label{sup4:lem:K}
For $p_1,p_2\in[0,3)$, $\delta>0$, $\xi\in(0,1]$ and $x,w\in\xi\Z^3$:
\[
\sum_{z\in\xi\Z^3}\xi^3(\xi+|x-z|)^{-p_1}(\xi+|z-w|)^{-p_2}e^{-\delta(|x-z|+|z-w|)}\le C_K\,h_{p_1+p_2-3}(\xi+|x-w|)\,e^{-\delta|x-w|/2}.
\]
\end{lemma}

\begin{proof}
Factor out $e^{-\delta|x-w|/2}$ by the triangle inequality and let $r\defeq\xi+|x-w|$. On $\{|x-z|\le|x-w|/2\}$ we have $\xi+|z-w|\ge r/2$, and $\sum_{|u|\le R}\xi^3(\xi+|u|)^{-p_1}\le C(R+\xi)^{3-p_1}$ since $p_1<3$; this region gives $\le Cr^{3-p_1-p_2}$. The region $\{|z-w|\le|x-w|/2\}$ is symmetric. On the rest, $|x-z|\le3|z-w|$ and the sum is $\le C\sum_{|u|\ge|x-w|/2}\xi^3(\xi+|u|)^{-p_1-p_2}e^{-\delta|u|/2}\le Ch_{p_1+p_2-3}(r)$. For $r\ge1$ the powers are absorbed by the remaining exponential.
\end{proof}

\begin{lemma}[cumulative minus free: the gain of the resolvent comparison, $d=3$]\label{sup4:lem:R}
Let $\breve G\defeq(-\Delta^\xi+\mu^2+bP_Y)^{-1}$ on $\xi\Z^3$, with $\mu^2\in[0,1]$, $b\in[b_-,b_+]\subset(0,\infty)$, $P_Y$ the projection onto constants on the blocks $Y$ of side $s_Y\in\{1,L\}$, and assume that $\breve G$ obeys \eqref{sup4:eq:Z1} (with $\xi$ in place of $\eta$). This holds for the rescaled cumulative propagators $\mathsf T_J^{-1}$ ($s_Y=1$) and $\mathsf T_{J+1}^{-1}$ ($s_Y=L$; the case $J=k'+1$ of \cref{sup4:cor:Z1} on the $\xi$-lattice, $\xi\Z^3$ being a union of $L$-blocks) of \cref{sup4:cor:Z1}, Step 1. Put $R\defeq\breve G-C^\xi$. Then
\begin{equation}\label{sup4:eq:R}
|\Delta_TD_SR(x_1,x_2)|\le C_R\,h_{|S|+\alpha|T|-1}(\xi+\rho)\,e^{-\delta_R\rho}\prod_{i\in T}|x_i-\tilde x_i|^\alpha .
\end{equation}
So $R$ and its H\"older quotients are bounded, one derivative costs at most a logarithm, and two derivatives (one per variable) cost $(\xi+\rho)^{-1}$: a gain of $2$ (up to the logarithm) over $C^\xi$.
\end{lemma}

\begin{proof}
$\breve G-C^\xi=\breve G\big((C^\xi)^{-1}-\breve G^{-1}\big)C^\xi=(1-\mu^2)\breve GC^\xi-b\,\breve GP_YC^\xi$; this is the identity of \cite[p.~437]{Ba3}.

\emph{Block part.} $P_Y=\sum_Y|Y|^{-1}\mathbf 1_Y\ip{\mathbf 1_Y}\cdot$, so the kernel of $\breve GP_YC^\xi$ is $\sum_Y|Y|^{-1}(\breve G\mathbf 1_Y)(x_1)(C^\xi\mathbf 1_Y)(x_2)$. $\breve G\mathbf 1_Y$ obeys the bounds of (A5): for $s_Y=1$ directly, and for $s_Y=L$ after rescaling by $L$ (the blocks become unit blocks), with constants depending on $L$. $C^\xi\mathbf 1_Y$ is a sum of $s_Y^3$ unit-block terms $C^\xi\mathbf 1_B$, $B\subset Y$, each bounded by \cref{sup4:lem:C1}. So this part is bounded by $Ce^{-\delta''\rho}\prod|\cdot|^\alpha\le Ch_\gamma(\xi+\rho)e^{-\delta''\rho}\prod|\cdot|^\alpha$ for every $\gamma$.

\emph{Convolution part.} The kernel is $\sum_z\xi^3\breve G(x_1,z)C^\xi(z,x_2)$; differences and quotients in $x_1$ fall on $\breve G$ and those in $x_2$ on $C^\xi$. By \eqref{sup4:eq:Z1} and \cref{sup4:lem:C} the factors are bounded by $(\xi+|\cdot|)^{-p_i}e^{-\delta|\cdot|}$ with $p_1=1+[1\in S]+\alpha[1\in T]<3$ and likewise $p_2$ (a quotient gives a minimum over two points, bounded by the sum of the two terms). \Cref{sup4:lem:K} gives $h_{p_1+p_2-3}=h_{|S|+\alpha|T|-1}$.
\end{proof}
